\par\medskip
\noindent
11.4. La descripción precedente de un acoplamiento (11.3.1), mediante
la expresión (11.2.1) para la componente $\ell$-primaria del grupo de
componentes conexas, solo es válida para $\ell\ne p$. Supongamos ahora
que $A_K$ tiene \underline{reducción semiestable} sobre $S$.
Definiremos de nuevo acoplamientos que son dualidades perfectas
\begin{equation}\tag{11.4.1}
\Phi_o(\ell)\otimes\Phi'_o(\ell)\longrightarrow\mathbb Q/\mathbb Z
\end{equation}
sin restricción alguna sobre $\ell$. Como en
\sisi{11.4}{11.3}, habría que establecer que se recupera (quizá salvo
un signo que convendría precisar) el acoplamiento definido en (1.2.1),
lo que demostraría la conjetura 1.3 en el caso de reducción
semiestable, como en el caso en que se restringe a las componentes
$\ell$-primarias para $\ell\ne p$.

Para definir (11.4.1), daremos esencialmente una expresión de
$\Phi_o(\ell)$ en términos del grupo de Barsotti--Tate $T_\ell(A_K)$

%%%%%%%%%%  scan idx 466  |  folio 455  %%%%%%%%%%

sobre $K$, que desempeñará en el presente caso el mismo papel que 11.2
en el caso tratado anteriormente.

\medskip
\noindent
\underline{Teorema 11.5.} \underline{Supongamos que $S$ es henseliano
y que $A_K$ tiene reducción semiestable sobre $S$. Sea $\ell$ un número
primo y consideremos el acoplamiento de monodromía} (9.1.2),
\underline{y, por tanto, el homomorfismo canónico}
\begin{equation*}
\widetilde u_\ell\sisi{;}{:}\underline M_\ell
 \longrightarrow\underline{\check M}'_\ell.
\end{equation*}
\underline{Pongamos
$\Phi(\ell)=\operatorname{coker}\widetilde u_\ell$. Entonces existe un
isomorfismo canónico}
\begin{equation}\tag{11.5.1}
\Phi_o(\ell)\simeq\Phi(\ell)\times_S s.
\end{equation}

De este enunciado se deduce inmediatamente la definición de un
acoplamiento perfecto (11.4.1). En efecto, aplicando 11.5 a $A'_K$ en
lugar de $A_K$ y utilizando (10.2.2), se obtiene también un isomorfismo
canónico
\begin{equation*}
\Phi'_o(\ell)\simeq\Phi'(\ell)\times_S s,
\qquad
\Phi'(\ell)=\operatorname{coker}\bigl(
 {}^t\!\widetilde u_\ell:\underline M'_\ell
 \longrightarrow\underline{\check M}_\ell\bigr).
\end{equation*}

\medskip
\noindent
\underline{Observaciones 11.5.2.} a) El teorema 11.5 implica, en
particular, que $\widetilde u_\ell$ es una isogenia, es decir, que
$u_\ell$ es separante, incluso en el caso $\ell=p$, donde esta
afirmación no es inmediata a partir de la definición de 9.6. Obsérvese
que no utilizamos aquí el teorema 10.4, considerablemente más profundo,
que se probará en la sección siguiente.

\smallskip
\noindent
b) Utilizando 10.4, que nos proporciona un acoplamiento de monodromía
\begin{equation*}
u:\underline M\otimes\underline M'\longrightarrow\mathbb Z_S,
\end{equation*}
y denotando por $\Phi$ el conúcleo del homomorfismo correspondiente
\begin{equation*}
\widetilde u:\underline M\longrightarrow\underline{\check M}',
\end{equation*}

%%%%%%%%%%  scan idx 467  |  folio 456  %%%%%%%%%%

obtenemos una formulación equivalente de 11.5, independiente de la
elección de $\ell$, mediante el isomorfismo canónico
\begin{equation}\tag{11.5.3}
\Phi_o\simeq\Phi\times_S s.
\end{equation}
